%% kapa: exam=1 round=31 year=2020 period=1 form=A subject=economics subject_ko=경제학원론 seq=29 number=69
%% source: Q-Net 감정평가사 기출문제 (cst00309) artlSeq=5207219 file=extracted/1008/5207219/A형/A형/2020년도 제31회 감정평가사 1차 시험_1교시_A형.pdf
%% conversion: kordoc 4.24.51 -> tools/gen1.py (auto); review=proofread-sonnet; flags=none
\begin{question}{69}
한국과 미국의 명목이자율은 각각 3\%, 2\% 이다. 미국의 물가상승률이 2\% 로 예상되며 현재 원/달러 환율은 1,000원일 때, 옳은 것을 모두 고른 것은? (단, 구매력평가설과 이자율평가설이 성립한다.)
\begin{framedtext}
ㄱ. 한국과 미국의 실질이자율은 같다. \par
ㄴ. 한국의 물가상승률은 3\%로 예상된다. \par
ㄷ. 원/달러 환율은 1,010원이 될 것으로 예상된다. \par
\end{framedtext}
\begin{choices}
  \item ㄱ
  \item ㄴ
  \item ㄱ, ㄴ
  \item ㄴ, ㄷ
  \item ㄱ, ㄴ, ㄷ
\end{choices}
\end{question}
\answer{5}
